\begin{definition}[Character wrappers]\label{mod:basic-charactertypes}
Let \(A\) be an additive monoid equipped with a topology.  A
\emph{continuous additive character} is a continuous homomorphism
\[
 \psi:(A,+)\longrightarrow\mathbf C^\times,
 \qquad
 \psi(0)=1,
 \quad
 \psi(x+y)=\psi(x)\psi(y).
\]
In the Lean formalization this is a continuous \texttt{AddChar} with
codomain \(\mathbf C^\times\).  For a monoid \(F\) equipped with a
topology, a \emph{continuous quasi-character} is a continuous monoid
homomorphism
\[
 \chi:F^\times\longrightarrow\mathbf C^\times.
\]
In the later local-field applications, \(F\) will be a field.

Continuous additive characters are closed under pointwise multiplication;
if \(A\) is an additive commutative group, they are also closed under
pointwise inverse.  Quasi-characters inherit the usual pointwise group
operations.  Every value is nonzero when viewed in \(\mathbf C\), since
the codomain is \(\mathbf C^\times\):
\[
 (\psi(x):\mathbf C)\ne0,
 \qquad
 (\chi(u):\mathbf C)\ne0.
\]
An additive character also has the explicit complex-valued realization
\[
 A\longrightarrow\mathbf C,
 \qquad
 x\longmapsto(\psi(x):\mathbf C),
\]
which is again an additive character with complex values and is continuous.
The Lean API exposes this realization as an explicit view, rather than as
a second implicit coercion.

\smallskip\noindent\emph{Generic pullback.}
For a continuous additive homomorphism \(f:A\to B\) and a continuous
additive character \(\psi\) of \(B\), define
\[
 f^*\psi=\psi\circ f.
\]
For a continuous homomorphism of unit groups
\(g:K^\times\to F^\times\) and a quasi-character \(\chi\) of \(F\),
define
\[
 g^*\chi=\chi\circ g.
\]
These pullbacks preserve the trivial character, pointwise multiplication,
inverse, and natural powers.  For additive characters the formulas are
\[
 \begin{aligned}
  f^*1&=1,&
  f^*(\psi\phi)&=(f^*\psi)(f^*\phi),\\
  f^*(\psi^{-1})&=(f^*\psi)^{-1},&
  f^*(\psi^n)&=(f^*\psi)^n,
 \end{aligned}
\]
and the same formulas hold for quasi-characters with \(g^*\) in place of
\(f^*\).  The inverse formula for additive characters requires the source
and target to be additive commutative groups.

\smallskip\noindent\emph{Trace and norm.}
Let \(R\) be a commutative topological ring and let \(S\) be a finite free
\(R\)-algebra, commutative for the trace construction, under the topology
hypotheses that make trace and norm continuous.  Pullback along trace and
norm gives, respectively,
\[
 \begin{aligned}
  \psi_S(x)
   &=\psi_R\!\left(\operatorname{Tr}_{S/R}(x)\right),
   &&x\in S,\\
  \chi_S(u)
   &=\chi_R\!\left(\operatorname{N}_{S/R}(u)\right),
   &&u\in S^\times.
 \end{aligned}
\]
The current generic Lean definitions use an
\texttt{IsModuleTopology R S} hypothesis to establish continuity of
trace and norm.  This topology witness is not part of the character types
themselves; the later local-field extension layer supplies the convenient
local-field specialization.
\LeanFile{LanglandsFirstMainLemma/Basic/CharacterTypes.lean}
\lean{LanglandsFirstMainLemma.ContinuousAddChar}
\leanok
\SourceText{Local notation, subsections Local fields, ideals, and characters and Norm and trace notation.}
\uses{mod:basic-phase}
\FormalizationContract
\end{definition}
